\documentclass[11pt,a4paper]{article}

% ── packages ──
\usepackage[utf8]{inputenc}
\usepackage[T1]{fontenc}
\usepackage{lmodern}
\usepackage[margin=2.5cm]{geometry}
\usepackage{amsmath,amssymb,bm}
\usepackage{hyperref}
\usepackage{listings}
\usepackage{xcolor}
\usepackage{graphicx}
\usepackage{booktabs}
\usepackage{enumitem}

% ── listing style ──
\definecolor{codebg}{RGB}{245,245,245}
\definecolor{codegreen}{RGB}{0,128,0}
\definecolor{codegray}{RGB}{100,100,100}
\definecolor{codepurple}{RGB}{128,0,128}
\lstdefinestyle{py}{
  language=Python,
  backgroundcolor=\color{codebg},
  basicstyle=\ttfamily\small,
  keywordstyle=\color{blue}\bfseries,
  stringstyle=\color{codegreen},
  commentstyle=\color{codegray}\itshape,
  emphstyle=\color{codepurple},
  frame=single,
  breaklines=true,
  showstringspaces=false,
  numbers=left,
  numberstyle=\tiny\color{codegray},
  tabsize=4,
  captionpos=b,
}
\lstset{style=py}

% ── shortcuts ──
\newcommand{\py}[1]{\lstinline[style=py]{#1}}
\newcommand{\file}[1]{\texttt{#1}}

% ══════════════════════════════════════════════════════════════════
\begin{document}

% ── title page ──
\begin{titlepage}
\centering
\vspace*{4cm}
{\Huge\bfseries Kalman Filter Experiment\\[6pt]
Codebase Explanation\par}
\vspace{1.5cm}
{\Large MLCOE Q2 --- Particle Flow Project\par}
\vspace{2cm}
{\large\today\par}
\vfill
{\itshape A self-contained guide covering architecture, execution flow,\\
TensorFlow design decisions, and the mathematical foundations\\
implemented in the \texttt{exp\_part1\_1a\_lgssm\_kf} pipeline.\par}
\end{titlepage}

\tableofcontents
\newpage

% ══════════════════════════════════════════════════════════════════
\section{Overview}
% ══════════════════════════════════════════════════════════════════

The Kalman Filter (KF) experiment is the simplest baseline in the
project: it applies the classical KF to a 2-D constant-velocity
Linear Gaussian State-Space Model (LGSSM).  Every subsequent
experiment (EKF, UKF, particle filters, particle flow) builds on the
same data-generation and metrics infrastructure, so understanding this
pipeline is essential.

\subsection{File Map}

\begin{table}[h]
\centering
\begin{tabular}{@{}ll@{}}
\toprule
\textbf{File} & \textbf{Role} \\
\midrule
\file{configs/ssm\_linear.yaml}            & Model \& noise parameters          \\
\file{src/models/ssm\_lgssm.py}            & \py{LGSSM} dataclass               \\
\file{src/data/generators.py}              & YAML parser + trajectory sampler   \\
\file{src/filters/kalman.py}               & \py{KalmanFilter} class            \\
\file{src/metrics/accuracy.py}             & RMSE, NEES, NLL metrics            \\
\file{src/experiments/exp\_part1\_1a\_lgssm\_kf.py} & Experiment driver         \\
\bottomrule
\end{tabular}
\caption{Files directly involved in the KF experiment.}
\end{table}

\subsection{Execution Flow}

\begin{enumerate}
\item \textbf{Parse configuration} ---
      \py{_parse_yaml_simple} reads \file{ssm\_linear.yaml}.
\item \textbf{Build model} ---
      \py{LGSSM.from_config} constructs matrices $A,B,C,D,Q,R,P_0$.
\item \textbf{Generate data} ---
      \py{LGSSM.sample(N)} draws a ground-truth trajectory
      $\{x_t\}_{t=1}^{N}$ and observations $\{y_t\}$.
\item \textbf{Filter} ---
      \py{KalmanFilter.filter(Y)} runs predict/update for every $t$.
\item \textbf{Evaluate} ---
      \py{compute_rmse} computes position and velocity RMSE.
\item \textbf{Save} --- CSV table + four-panel PNG figure.
\end{enumerate}


% ══════════════════════════════════════════════════════════════════
\section{The State-Space Model}
\label{sec:ssm}
% ══════════════════════════════════════════════════════════════════

\subsection{Mathematical Formulation}

The discrete-time LGSSM is
\begin{align}
  \bm{x}_t &= A\,\bm{x}_{t-1} + B\,\bm{v}_t,
    \qquad \bm{v}_t \sim \mathcal{N}(\bm{0},\,I), \label{eq:trans}\\
  \bm{y}_t &= C\,\bm{x}_t + D\,\bm{w}_t,
    \qquad \bm{w}_t \sim \mathcal{N}(\bm{0},\,I), \label{eq:obs}
\end{align}
with process-noise covariance $Q = BB^\top$ and measurement-noise
covariance $R = DD^\top$.

\subsection{Constant-Velocity Model (ssm\_linear.yaml)}

The state vector is $\bm{x} = [p_x,\, v_x,\, p_y,\, v_y]^\top$
with $\Delta t = 1$:
\[
A = \begin{bmatrix}1&1&0&0\\0&1&0&0\\0&0&1&1\\0&0&0&1\end{bmatrix},
\quad
B_{\text{raw}} = \begin{bmatrix}0.5&0\\1&0\\0&0.5\\0&1\end{bmatrix},
\quad
C = \begin{bmatrix}1&0&0&0\\0&0&1&0\end{bmatrix}.
\]
The YAML specifies $\sigma_a=0.5$ (acceleration noise standard
deviation) and $\sigma_z=5.0$ (measurement noise standard deviation).
$B$ and $D$ are scaled in code:
$B = \sigma_a B_{\text{raw}}$, $D = \sigma_z D_{\text{raw}}$.

\subsection{Initial Covariance ($P_0$)}

The config sets
$P_0 = \operatorname{diag}(10,\,5,\,10,\,5)$
and $\bm{m}_0 = \bm{0}$.  This is a diagonal prior expressing
moderate uncertainty in position (10) and velocity (5).

\subsection{LGSSM Class (\file{ssm\_lgssm.py})}

\py{LGSSM} is a Python \py{@dataclass} holding the eight tensors
$(A,B,C,D,\bm{m}_0,P_0,Q,R)$ plus integer dimensions.
Its two methods are:

\begin{description}[leftmargin=1.5cm]
\item[\py{from_config}] Static factory that parses the config dict,
  scales $B$/$D$ by noise standard deviations, and computes $Q$/$R$ via
  \py{tf.matmul(B, B, transpose_b=True)}.
\item[\py{sample}] Generates a trajectory by iterating
  Eqs.~\eqref{eq:trans}--\eqref{eq:obs}.  Uses
  \py{tfd.MultivariateNormalTriL} for the initial draw and
  \py{tfd.MultivariateNormalDiag} for i.i.d.\ noise.
\end{description}

\paragraph{TensorFlow detail.}
The Cholesky decomposition $P_0 = LL^\top$ is computed via
\py{tf.linalg.cholesky(P0)} and passed to
\py{MultivariateNormalTriL(loc=m0, scale_tril=L)}.  TFP's TriL
parameterisation is preferred because it avoids redundant
eigendecompositions that \py{MultivariateNormalFullCovariance} would
perform internally.


% ══════════════════════════════════════════════════════════════════
\section{Data Generation (\file{generators.py})}
% ══════════════════════════════════════════════════════════════════

\subsection{\py{_parse_yaml_simple}}

A dependency-free YAML parser (no PyYAML required).  It handles:
\begin{itemize}
\item Scalars (\py{seed: 42}),
\item Indented sub-dicts (\py{dimensions}, \py{params}),
\item Row-major matrices (\py{A}, \py{B_raw}, \py{C}, \py{D}),
\item Diagonal vectors (\py{Sigma0_diag}).
\end{itemize}

\subsection{\py{generate_lgssm_from_yaml}}

Orchestrates the full data pipeline:

\begin{lstlisting}[caption={Data generation entry point.}]
model, X_true, Y_obs, data_dict = generate_lgssm_from_yaml(config_path)
\end{lstlisting}

Returns the \py{LGSSM} instance, the true state matrix
$X \in \mathbb{R}^{N \times 4}$, the observation matrix
$Y \in \mathbb{R}^{N \times 2}$, and a convenience dictionary
splitting each component into named tensors for plotting.


% ══════════════════════════════════════════════════════════════════
\section{The Kalman Filter (\file{kalman.py})}
\label{sec:kf}
% ══════════════════════════════════════════════════════════════════

\subsection{Constructor}

\begin{lstlisting}[caption={KalmanFilter initialisation.}]
kf = KalmanFilter(A=model.A, C=model.C,
                  x0=x0_reshaped, P0=model.P0,
                  Q=model.Q, R=model.R)
\end{lstlisting}

Key design decisions:
\begin{itemize}
\item \textbf{Column-vector state}: $\bm{x}$ is stored as shape
  $(n, 1)$ so that all updates use \py{tf.matmul} without manual
  transpose.
\item \textbf{Flexible noise input}: $Q$ and $R$ can be supplied
  directly or derived from $B$/$D$ via $Q = BB^\top$.
  If neither is given, a zero matrix is used (noiseless system).
\end{itemize}

\subsection{Predict Step}
\label{sec:predict}

The time update propagates the state estimate and covariance one step
forward:
\begin{align}
  \hat{\bm{x}}_{t|t-1} &= A\,\hat{\bm{x}}_{t-1|t-1},
    \label{eq:kf-pred-mean}\\
  P_{t|t-1} &= A\,P_{t-1|t-1}\,A^\top + Q.
    \label{eq:kf-pred-cov}
\end{align}

\begin{lstlisting}[caption={Predict step implementation.}]
def predict(self):
    self.x = tf.matmul(self.A, self.x)
    self.P = tf.matmul(self.A,
                tf.matmul(self.P, self.A, transpose_b=True)) + self.Q
    return self.x, self.P
\end{lstlisting}

\paragraph{TensorFlow note.}
\py{tf.matmul(P, A, transpose_b=True)} computes $PA^\top$ in a
single fused BLAS call without creating the transposed matrix
explicitly, saving memory and one kernel launch.

\subsection{Update Step}
\label{sec:update}

Given a new observation $\bm{z}_t$:
\begin{align}
  \bm{\nu}_t &= \bm{z}_t - C\,\hat{\bm{x}}_{t|t-1},
    &\text{(innovation)} \label{eq:innov}\\
  S_t &= C\,P_{t|t-1}\,C^\top + R,
    &\text{(innovation covariance)} \label{eq:S}\\
  K_t &= P_{t|t-1}\,C^\top\,S_t^{-1},
    &\text{(Kalman gain)} \label{eq:gain}\\
  \hat{\bm{x}}_{t|t} &= \hat{\bm{x}}_{t|t-1} + K_t\,\bm{\nu}_t,
    &\text{(filtered mean)} \label{eq:xfilt}\\
  P_{t|t} &= (I - K_t C)\,P_{t|t-1}.
    &\text{(Riccati update)} \label{eq:Pfilt}
\end{align}

\subsubsection{Kalman Gain via Linear Solve}

The gain is \emph{not} computed using an explicit matrix inverse.
Instead, the code solves a linear system:

\begin{lstlisting}[caption={Gain computation via \py{tf.linalg.solve}.}]
CP = tf.matmul(self.C, self.P)        # (m, n)
K_T = tf.linalg.solve(S, CP)          # S @ K^T = CP
K = tf.transpose(K_T)                 # (n, m)
\end{lstlisting}

\py{tf.linalg.solve} internally calls LAPACK's \texttt{dgesv}
(LU with partial pivoting), which is both faster and more numerically
stable than forming $S^{-1}$ via \py{tf.linalg.inv}.  The latter
computes all columns of the inverse even though only $CP$ is needed.

\subsubsection{Joseph-Stabilised Update (Optional)}

When \py{joseph=True}, the covariance update uses the
symmetry-preserving Joseph form:
\begin{equation}
  P_{t|t} = (I - K_tC)\,P_{t|t-1}\,(I - K_tC)^\top
           + K_t\,R\,K_t^\top.
  \label{eq:joseph}
\end{equation}
This is algebraically identical to the Riccati form
Eq.~\eqref{eq:Pfilt} but better preserves positive-definiteness
under floating-point rounding.

\begin{lstlisting}[caption={Joseph form implementation.}]
if joseph:
    IKC = I - tf.matmul(K, self.C)
    self.P = tf.matmul(IKC,
                tf.matmul(self.P, IKC, transpose_b=True)) \
           + tf.matmul(K, tf.matmul(self.R, K, transpose_b=True))
else:
    self.P = tf.matmul(I - tf.matmul(K, self.C), self.P)
\end{lstlisting}

\subsection{The \py{filter} Method}

Iterates predict--update over all $N$ measurements:

\begin{lstlisting}[caption={Full filtering loop.}]
def filter(self, Z, *, joseph=False, use_solve=True):
    for t in range(N):
        xp, Pp = self.predict()
        xf, Pf, *_ = self.update(Z[t], joseph=joseph,
                                  use_solve=use_solve)
    return {"x_pred": ..., "P_pred": ...,
            "x_filt": ..., "P_filt": ...}
\end{lstlisting}

Returns stacked tensors of shape $(N, n, 1)$ for means and
$(N, n, n)$ for covariances, enabling vectorised post-processing.


% ══════════════════════════════════════════════════════════════════
\section{Metrics (\file{accuracy.py})}
% ══════════════════════════════════════════════════════════════════

Only \py{compute_rmse} is used by the KF experiment.  The other
functions (NEES, NIS, NLL, innovation whiteness) are consumed by the
EKF/UKF experiments downstream.

\subsection{\py{compute_rmse}}

\begin{lstlisting}[caption={RMSE implementation.}]
def compute_rmse(estimates, ground_truth):
    if estimates.shape.ndims == 3:
        estimates = tf.squeeze(estimates, axis=-1)
    return float(tf.sqrt(tf.reduce_mean(
               (estimates - ground_truth) ** 2)))
\end{lstlisting}

\begin{equation}
  \text{RMSE} = \sqrt{\frac{1}{N \cdot d}\sum_{t=1}^{N}
    \|\hat{\bm{x}}_t - \bm{x}_t\|^2}.
  \label{eq:rmse}
\end{equation}

The experiment computes two separate RMSE values:
\begin{itemize}
\item \textbf{Position RMSE}: using dimensions $(p_x, p_y)$,
      i.e.\ columns 0 and 2.
\item \textbf{Velocity RMSE}: using dimensions $(v_x, v_y)$,
      i.e.\ columns 1 and 3.
\end{itemize}

\paragraph{TensorFlow detail.}
\py{tf.squeeze(estimates, axis=-1)} removes the trailing
singleton dimension $(N, n, 1) \to (N, n)$ so that element-wise
subtraction broadcasts correctly against the $(N, n)$ ground truth.


% ══════════════════════════════════════════════════════════════════
\section{Experiment Script}
\label{sec:experiment}
% ══════════════════════════════════════════════════════════════════

\file{exp\_part1\_1a\_lgssm\_kf.py} is a self-contained driver with
a single public function and a CLI entry point.

\subsection{\py{run_experiment(config_path, out_dir)}}

\begin{enumerate}
\item \textbf{Load \& generate}: calls \py{generate_lgssm_from_yaml}.
\item \textbf{Reshape $\bm{m}_0$}: ensures the initial state is
  $(n, 1)$.  A try/except handles both static and dynamic TF shapes:

\begin{lstlisting}[caption={Shape handling for dynamic graphs.}]
try:
    if model.m0.shape.ndims == 1:
        x0_reshaped = tf.reshape(model.m0, [-1, 1])
    else:
        x0_reshaped = model.m0
except (AttributeError, ValueError):
    x0_rank = tf.rank(model.m0)
    x0_reshaped = tf.cond(tf.equal(x0_rank, 1),
                   lambda: tf.reshape(model.m0, [-1, 1]),
                   lambda: model.m0)
\end{lstlisting}

The first branch uses \py{.shape.ndims} (a Python integer available
in eager mode).  The fallback uses \py{tf.rank} and \py{tf.cond},
which work inside \py{@tf.function} graphs where static shape info
may be unavailable.

\item \textbf{Filter}: \py{kf.filter(Y_obs)}.
\item \textbf{Extract components}:

\begin{lstlisting}[caption={Splitting position and velocity.}]
xf = tf.squeeze(out["x_filt"], axis=-1)   # (N, 4)
xf_pos = tf.stack([xf[:, 0], xf[:, 2]], axis=1)  # (N, 2)
xf_vel = tf.stack([xf[:, 1], xf[:, 3]], axis=1)  # (N, 2)
\end{lstlisting}

\item \textbf{Compute RMSE}: separately for position and velocity.
\item \textbf{Save CSV}: one row per timestep with all true, observed,
  and filtered values.
\item \textbf{Plot}: four-panel matplotlib figure (position $x$,
  position $y$, velocity $v_x$, velocity $v_y$).
\end{enumerate}

\subsection{CLI Interface}

\begin{lstlisting}[language=bash,style=py]
python -m src.experiments.exp_part1_1a_lgssm_kf \
    --config configs/ssm_linear.yaml \
    --out_dir reports/1_LinearGaussianSSM/figures
\end{lstlisting}


% ══════════════════════════════════════════════════════════════════
\section{TensorFlow API Summary}
% ══════════════════════════════════════════════════════════════════

Table~\ref{tab:tf} lists every TensorFlow function used in the KF
pipeline and explains why it was chosen.

\begin{table}[h]
\centering
\small
\begin{tabular}{@{}p{4.8cm}p{8.2cm}@{}}
\toprule
\textbf{TF Function} & \textbf{Role \& Rationale} \\
\midrule
\py{tf.matmul(A, B, transpose_b=True)}
  & Fused $AB^\top$; avoids allocating the transpose. \\
\py{tf.linalg.solve(S, CP)}
  & Computes $S^{-1}(CP)$ via LU factorisation;
    numerically superior to explicit \py{tf.linalg.inv}. \\
\py{tf.linalg.inv(S)}
  & Fallback gain path (disabled by default). \\
\py{tf.linalg.cholesky(P0)}
  & Lower-triangular factor for sampling
    $\bm{x}_0 \sim \mathcal{N}(\bm{m}_0, P_0)$. \\
\py{tf.linalg.diag(v)}
  & Builds diagonal $P_0$ and $Q$/$R$ from config vectors. \\
\py{tf.eye(n)}
  & Identity matrix for $(I - KC)P$ and Joseph form. \\
\py{tf.reshape(x, [-1, 1])}
  & Ensures column-vector shape for \py{matmul}. \\
\py{tf.squeeze(x, axis=-1)}
  & Removes trailing singleton $(N,n,1) \to (N,n)$. \\
\py{tf.stack(list, axis=0)}
  & Stacks per-timestep results into batch tensor. \\
\py{tf.convert_to_tensor}
  & Normalises numpy/list inputs to TF tensors. \\
\py{tf.rank}, \py{tf.cond}
  & Graph-safe dynamic shape dispatch. \\
\py{tf.reduce_mean}, \py{tf.sqrt}
  & RMSE computation. \\
\py{tfd.MultivariateNormalTriL}
  & Cholesky-parameterised Gaussian for initial state. \\
\py{tfd.MultivariateNormalDiag}
  & Diagonal Gaussian for i.i.d.\ process/measurement noise. \\
\bottomrule
\end{tabular}
\caption{TensorFlow functions and their roles in the KF pipeline.}
\label{tab:tf}
\end{table}


% ══════════════════════════════════════════════════════════════════
\section{Interaction Diagram}
% ══════════════════════════════════════════════════════════════════

The call graph for a single experiment run:

\begin{verbatim}
exp_part1_1a_lgssm_kf.run_experiment()
  |
  +-> generators.generate_lgssm_from_yaml(config_path)
  |     +-> generators._parse_yaml_simple(config_path)
  |     +-> LGSSM.from_config(cfg)      # builds A,B,C,D,Q,R,P0
  |     +-> LGSSM.sample(N, seed)       # generates X_true, Y_obs
  |
  +-> KalmanFilter.__init__(A, C, x0, P0, Q, R)
  |
  +-> KalmanFilter.filter(Y_obs)
  |     +-- for t in range(N):
  |           +-> KalmanFilter.predict()    # Eqs. (3)-(4)
  |           +-> KalmanFilter.update(z_t)  # Eqs. (5)-(9)
  |
  +-> accuracy.compute_rmse(xf_pos, X_true_pos)
  +-> accuracy.compute_rmse(xf_vel, X_true_vel)
  |
  +-> matplotlib: save 4-panel figure
  +-> CSV: save tabular results
\end{verbatim}


% ══════════════════════════════════════════════════════════════════
\section{Design Decisions \& Potential Improvements}
% ══════════════════════════════════════════════════════════════════

\begin{enumerate}
\item \textbf{Eager-mode loop.}
  The \py{for t in range(N)} loop in \py{filter()} runs in eager
  mode.  For $N=200$ this is negligible, but for longer sequences a
  \py{tf.while_loop} or \py{tf.scan} would enable XLA compilation.

\item \textbf{$P_0$ initialisation.}
  The current diagonal $P_0$ is hand-tuned.  Alternatives include:
  \begin{itemize}
  \item \emph{Diffuse prior}: $P_0 = \alpha I$ with large $\alpha$.
  \item \emph{Steady-state DARE}: solve the Discrete Algebraic
        Riccati Equation for the asymptotic $P$.
  \item \emph{Physics-informed}:
        $\sigma_{p} = \sigma_z$,
        $\sigma_{v} = \sigma_z / \Delta t$.
  \end{itemize}

\item \textbf{Joseph form.}
  The default Riccati update is fine for double precision but can
  lose positive-definiteness in single precision (\py{float32}).
  Setting \py{joseph=True} prevents this at the cost of one extra
  matrix multiply.

\item \textbf{Linear solve vs inverse.}
  The \py{use_solve=True} default is the correct modern choice.
  The \py{tf.linalg.inv} path exists only as a pedagogical reference
  and should not be used in production.
\end{enumerate}


% ══════════════════════════════════════════════════════════════════
\section{Conclusion}
% ══════════════════════════════════════════════════════════════════

This document traced the complete execution path of the KF experiment
from YAML configuration to CSV/PNG output.  The pipeline is
intentionally simple---six files, no training loop, no gradient
tape---because the KF serves as the \emph{analytical ground truth}
for the linear-Gaussian case.  Every subsequent filter in the project
(EKF, UKF, bootstrap PF, LEDH, EDH, Dai--Daum) is benchmarked
against this baseline, making correctness here critical.

\end{document}
